\section{Introduction}
\label{sec:introduction}

Modern passage-ranking systems are commonly deployed as multi-stage
pipelines: a cheap first-stage retriever, often BM25, retrieves a
candidate pool, and a more expensive neural model reranks the top
candidates~\citep{robertson2009probabilistic,
nogueira2019passage,nogueira2020monot5,yates2021pretrained}.
This setting underlies standard passage-ranking evaluations such as
MS MARCO, the TREC Deep Learning Tracks, and BEIR, where rerankers are
measured by their ability to improve a first-stage ranking under graded
relevance metrics such as nDCG@10~\citep{bajaj2016msmarco,
craswell2020trecDL,craswell2021trecDL,thakur2021beir,
jarvelin2002cumulated}. Recent large language model (LLM) rerankers
revisit this pipeline in a zero-shot form: instead of training a
task-specific cross-encoder, the model is prompted to judge, compare, or
order retrieved passages~\citep{sachan2022upr,sun2023rankgpt,
ma2023zeroshotlistwise,qin2024prp,zhuang2024setwise,
pradeep2023rankvicuna,pradeep2023rankzephyr}.

LLM zero-shot reranking has largely been organized around four prompting
paradigms. Pointwise methods estimate relevance for each query--document
pair independently, either with supervised sequence-to-sequence relevance
labels or zero-shot likelihood-style scoring~\citep{nogueira2020monot5,
sachan2022upr}. Pairwise methods ask the model to choose between two
documents and aggregate local preferences into a ranking, often with
additional procedures to reduce order bias or comparison
cost~\citep{pradeep2021expandomonoduo,qin2024prp}. Listwise methods
present a short list of documents and ask the model to emit a permutation
or ordered list of document identifiers~\citep{ma2023zeroshotlistwise,
sun2023rankgpt,pradeep2023rankvicuna,pradeep2023rankzephyr,
yoon2024listt5}. Setwise methods occupy a useful middle ground: one LLM
call compares more than two documents, but the model returns only a short
label, such as the identifier of the most relevant document, rather than
a full permutation~\citep{zhuang2024setwise}. This smaller output
interface makes setwise prompting easier to parse than listwise
permutation generation while allowing each decision to condition on more
context than pairwise comparison.

Existing setwise rerankers are usually designed around small local
windows. A typical call sees a query and a window of $c{+}1$ candidate
passages, selects one extreme item, and embeds that local decision inside
a sorting procedure such as heapsort or bubblesort
~\citep{qin2024prp,zhuang2024setwise}. This design was natural when
context length was a binding constraint: the prompt stayed small, the
output was easy to validate, and the sorting algorithm supplied the
global ranking structure. The same locality appears in related listwise
reranking systems, where sliding-window, tournament-style, or
cache-assisted procedures are used to avoid asking the model to rank too
many passages at once~\citep{sun2023rankgpt,ma2023zeroshotlistwise,
yoon2024listt5}. As a result, small-window setwise reranking has become a
strong and practical zero-shot baseline. It also inherits a question from
the original context-limited setting: are local windows still necessary
when the whole reranking pool fits in context?

The context-length regime has changed for many open-weight rerankers.
For the common first-stage reranking setting with $N \le 100$, a query
and all candidate passages can now often fit in a single prompt. With
passage length $\ell=512$, a query and $50$ passages require roughly
$26$K input tokens before output reserve, below the advertised or
documented long-context budgets of the Qwen3.5, Llama-3.1, and
Ministral-3 families considered in this paper~\citep{qwen2026qwen35,
dubey2024llama3,liu2026ministral3}. This observation does not make
small-window setwise reranking obsolete. Long prompts can increase
input-token cost and per-call latency, and long-context models do not
necessarily use all prompt positions uniformly~\citep{liu2024lost}.
Moreover, label-based prompting can be sensitive to option identifiers,
candidate order, and presentation position~\citep{zheng2024mcqbias,
shi2025positionbias,pradeep2023rankzephyr}. The changed context regime
therefore motivates a missing baseline rather than a replacement claim:
if the current live pool fits in context, how much of the usual setwise
sorting machinery is still needed?

We study this question through \textbf{MaxContext}, a zero-shot setwise
reranking family in which every LLM call observes the current whole live
pool rather than a local candidate window. MaxContext keeps the setwise
output interface deliberately small. Documents are labeled numerically
$1,\ldots,|L|$, and the model returns either one extreme document or both
extreme documents from the live pool. \textsc{MaxContext-TopDown} asks
for the best document and places it at the top of the ranking.
\textsc{MaxContext-BottomUp} asks for the worst document and places it at
the bottom. \textsc{MaxContext-DualEnd} asks for the best and worst
documents jointly, placing one at each end and removing two documents
from the live pool in a single call. The single-extreme variants close
the final two-document residual with a deterministic BM25 endgame.

This three-way design separates two questions that are easy to conflate.
The first is whether whole-pool setwise selection is a viable alternative
to small-window setwise sorting. The second is whether eliciting both
ends of the ranking in the same call preserves effectiveness while
reducing the number of serial LLM calls. The call-count comparison within
MaxContext is deterministic: for a pool of size $N$,
\textsc{MaxContext-DualEnd} uses $\lfloor N/2 \rfloor$ LLM calls,
whereas either single-extreme MaxContext variant uses $N{-}2$ LLM calls
plus the BM25 two-document endgame. At $N{=}100$, this is $50$ versus
$98$ LLM calls. We therefore evaluate whether the joint best--worst
interface delivers this call reduction without a corresponding loss in
nDCG@10.

Our experiments compare seven methods: standard TopDown setwise
reranking with bubblesort and heapsort, standard BottomUp setwise
reranking with bubblesort and heapsort, and the three MaxContext
variants. To test whether the effect is specific to one chat template,
tokenizer, or long-context implementation, we evaluate nine open-weight
models from the Qwen3.5, Llama-3.1, and Ministral-3 families. The main
matrix covers TREC DL19/20 and six BEIR domains, with pool sizes
$N \in \{10,20,30,40,50,100\}$. We measure effectiveness with per-query
nDCG@10~\citep{jarvelin2002cumulated}, assess significance using paired
bootstrap tests~\citep{efron1993bootstrap,smucker2007comparison,
urbano2019significance}, and apply Bonferroni correction across the
seven methods~\citep{bonferroni1936teoria}. We also run an input-order
stability protocol on DL19 with three representative models and ten
randomized repetitions, motivated by known order and position effects in
long-context and label-based LLM prompting~\citep{liu2024lost,
zheng2024mcqbias,shi2025positionbias,pradeep2023rankzephyr}.

We frame the efficiency claim carefully. MaxContext-DualEnd is designed
to reduce the number of serial LLM calls, and we report wall-clock time
to measure the practical effect of that reduction. We also report input,
output, and total token counts, but we do not claim token efficiency:
whole-pool prompts can consume more input tokens than small-window
setwise calls even when they require fewer calls. This distinction is
important because latency, forward passes, input tokens, output tokens,
and FLOPs capture different aspects of the efficiency--effectiveness
tradeoff for LLM rerankers~\citep{peng2025flopsreranking}. Our central
empirical question is therefore whether whole-pool dual-end elicitation
offers a better call-latency tradeoff while maintaining reranking
effectiveness.

Our contributions are threefold. First, we introduce MaxContext, a
whole-pool setwise reranking family with best-only, worst-only, and joint
best--worst output policies, yielding a deterministic call-count gap of
$\lfloor N/2 \rfloor$ versus $N{-}2$ within the MaxContext family.
Second, we provide a controlled comparison against standard setwise
TopDown and BottomUp rerankers instantiated with bubblesort and heapsort,
covering seven methods across nine models, eight datasets, and six pool
sizes. Third, we evaluate both effectiveness and stability, using
per-query nDCG@10, paired-bootstrap significance testing with Bonferroni
correction, wall-clock and token accounting, and an explicit input-order
stability protocol.